% BHU PYQ -- Manually transcribed & error-corrected
\documentclass[11pt,a4paper]{article}

\usepackage[a4paper,top=2.5cm,bottom=2.5cm,left=2.8cm,right=2.8cm,headheight=14pt]{geometry}
\usepackage{amsmath,amssymb}
\usepackage[T1]{fontenc}
\usepackage[utf8]{inputenc}
\usepackage{lmodern}
\usepackage{microtype}
\usepackage{fancyhdr}
\usepackage{enumitem}
\usepackage{hyperref}
\usepackage{lastpage}
\usepackage{parskip}

\hypersetup{colorlinks=true,urlcolor=black,linkcolor=black,
  pdftitle={BPT-504: Electronic Devices and Circuits -- BHU PYQ},pdfauthor={Banaras Hindu University}}

\pagestyle{fancy}\fancyhf{}
\renewcommand{\headrulewidth}{0.4pt}\renewcommand{\footrulewidth}{0.4pt}
\fancyhead[L]{\small   \textit{Banaras Hindu University}}
\fancyhead[R]{\small   \textit{BPT-504: Electronic Devices and Circuits}}
\fancyfoot[C]{\small Page  \thepage\ of \pageref{LastPage}}


\newlist{parts}{enumerate}{2}
\setlist[parts,1]{label=(\alph*),leftmargin=*,itemsep=5pt,topsep=3pt}
\setlist[parts,2]{label=(\roman*),leftmargin=*,itemsep=3pt,topsep=2pt}

\newcommand{\pts}[1]{\hfill   \textit{[#1]}}


\begin{document}


\begin{center}
  {\large\bfseries BANARAS HINDU UNIVERSITY}\\[2pt]
  {
\normalsize B.Sc. (Hons.) Semester V Examination, 2013-14}\\[6pt]
  \rule{0.8\linewidth}{0.5pt}\\[6pt]
  {\large\bfseries Paper No. BPT-504: Electronic Devices and Circuits}\\[4pt]
  {
\normalsize Time: 3 Hours\hfill Maximum Marks: 70}
\end{center}
\rule{\linewidth}{0.4pt}
\smallskip



\noindent   \textit{Instructions: Answer all questions. Symbols have their usual meanings.}
\bigskip




\noindent   \textbf{Question 1.} Answer in brief any   \textbf{five} of the following: \pts{5 $ \times$ 4 = 20}

\begin{parts}
  \item What do you mean by the intrinsic standoff ratio ($\eta$) of a UJT? If the interbase resistance $R_{BB} = 8\, \text{k}\Omega$ and the base-1 resistance $R_{B1} = 4.8\, \text{k}\Omega$, find the value of $\eta$.
  \item An amplifier with a voltage gain of $A_v = -1000$ has a gain change of $20\%$ due to temperature change. Calculate the percentage change in gain of the feedback amplifier if a negative feedback fraction of $\beta = -0.1$ is incorporated.
  \item Differentiate among class A, class B, and class C amplifiers on the basis of their operating points and maximum theoretical efficiencies.
  \item What is the main difference between oscillators and multivibrators? Draw the circuit diagram of a BJT-based astable multivibrator.
  \item Give the truth table and Boolean expression for a XOR gate and implement it using NAND gates only.
  \item Draw the circuit diagram of an envelope detector demodulator for an amplitude-modulated (AM) wave and explain its working qualitatively.
\end{parts}

\medskip\rule{\linewidth}{0.2pt}




\noindent   \textbf{Question 2.}

\begin{parts}
  \item Explain the construction, working, and characteristic curves of an enhancement-type MOSFET. How does it differ from a JFET? What is the significance of the threshold voltage ($V_T$) in an enhancement mode MOSFET? \pts{7}
  \item Draw the circuit diagram of a non-inverting amplifier using an OP-AMP and derive the expression for its closed-loop voltage gain. How will you modify this circuit to act as a voltage follower? \pts{7}
\end{parts}

\medskip\rule{\linewidth}{0.2pt}




\noindent   \textbf{Question 3.}

\begin{parts}
  \item Draw the hybrid model (h-parameter) equivalent circuit of a BJT in common-emitter (CE) configuration. Derive the expressions for its voltage gain, current gain, input impedance, and output admittance. \pts{7}
  \item Explain the working of a Silicon Controlled Rectifier (SCR) with the help of a two-transistor equivalent circuit. Explain why it is preferred to use silicon instead of germanium to construct a controlled rectifier. \pts{7}
\end{parts}
\hfill   \textbf{OR}

\begin{parts}
  \item Explain in brief the high-frequency hybrid-$\pi$ model of a CE BJT amplifier and discuss its advantages. \pts{5}
  \item Refer to the basic common-base BJT amplifier configuration. The h-parameter values for the transistor are:
    \[ h_{ib} = 40\,\Omega, \quad h_{rb} = 3  \times 10^{-4}, \quad h_{fb} = -0.99, \quad h_{ob} = 0.5\,\mu \text{mhos} \]
    Determine: (i) the current gain taking source resistance $R_s = 600\,\Omega$ and load resistance $R_L = 10\, \text{k}\Omega$ into account, and (ii) the input resistance. \pts{5}
  \item Give the construction and working of a Unijunction Transistor (UJT). Explain the condition under which a UJT exhibits negative resistance. \pts{4}
\end{parts}

\medskip\rule{\linewidth}{0.2pt}




\noindent   \textbf{Question 4.}

\begin{parts}
  \item Draw the circuit diagram of a common-emitter BJT amplifier using RC coupling. Derive the equations to show that the gain of this amplifier decreases at the low-frequency range due to coupling and bypass capacitors. \pts{5}
  \item Give the truth table of a full adder. Obtain the Boolean expressions for sum ($S$) and carry ($C_o$) from this table in Sum-of-Products (SOP) form. Simplify these expressions using a Karnaugh map (K-map) and implement the circuit using basic logic gates. \pts{5}
  \item The maximum and minimum amplitudes of an amplitude-modulated (AM) wave are $10\, \text{V}$ and $2\, \text{V}$ respectively. Calculate the carrier amplitude and the modulation index. \pts{4}
\end{parts}

\medskip\rule{\linewidth}{0.2pt}




\noindent   \textbf{Question 5.}

\begin{parts}
  \item Explain the working of a class B push-pull amplifier with a suitable circuit diagram. Show mathematically how a push-pull amplifier reduces even-harmonic distortions. \pts{7}
  \item What is meant by negative feedback in amplifiers? Find the closed-loop voltage gain, input resistance, and output resistance of a voltage-series feedback amplifier. Draw its equivalent BJT circuit diagram. \pts{7}
\end{parts}
\hfill   \textbf{OR}

\begin{parts}
  \item Assume that it is desired to design a Colpitts oscillator so that its oscillation frequency is $4\, \text{MHz}$. If the inductor is $L = 50\,\mu \text{H}$ and the desired feedback fraction is $\beta = 1/20$, calculate:
    
\begin{parts}
      \item The required value of equivalent tank capacitance. \pts{3}
      \item The required values for capacitances $C_1$ and $C_2$. \pts{2}
      \item The minimum voltage gain $A_v$ required for the CE amplifier to sustain oscillations. \pts{2}
    \end{parts}
  \item Draw the basic RC phase shift oscillator circuit and explain its working. As per the Barkhausen criterion, what should be the minimum amplifier gain for sustained oscillations? What will be the frequency of oscillation? How are the loading effects of this circuit eliminated in a Wien bridge oscillator? \pts{7}
\end{parts}

\vfill
\rule{\linewidth}{0.4pt}

\begin{center}\small
  \href{https://bhu-exam.vercel.app/}{Click here to visit our website}\\[4pt]
  \href{https://me-aryan.vercel.app/}{Click here to visit our contributor}\\[4pt]
  \href{https://quashan-me.vercel.app/}{Click here to visit our contributor}
\end{center}

\end{document}
